\section{拉格朗日道路的马斯洛夫指标}
\label{sec:app-c-maslov-paths}

本节逐字依循\personRobbin{}与\personSalamon{}在 (RS95) 第 5 节中的方法（仅作少量记号调整），概述与每一
对拉格朗日道路 $(\Lambda,\Lambda')$ 相伴的马斯洛夫指标之基本性质。该方法又以他们在
(RS93) 中对马斯洛夫指标的详细叙述为基础。
\footnote{译者注：原书此处把 Salamon 误拼成 ``Salmon''；译文采用受控人名“萨拉蒙”。}

设 $(S,\Omega)$ 是辛向量空间，以 $\mathcal L(S)$ 表示其中拉格朗日子空间的集合。
(RS93) 所定义的马斯洛夫指标，为每一对拉格朗日道路
$\Lambda,\Lambda'\colon[a,b]\to\mathcal L(S)$ 指定一个半整数
$\mu(\Lambda,\Lambda')$。它由下列公理刻画。

\begin{enumerate}[label=(\arabic*)]
\item \stfhPropertyLabel{（自然性）}对任意辛微分同胚
$\Psi\colon(S,\Omega)\to(S',\Omega')$，
\[
 \mu(\Psi\Lambda,\Psi\Lambda')=\mu(\Lambda,\Lambda').
\]

\item \stfhPropertyLabel{（同伦性）}马斯洛夫指标在端点固定的同伦下不变。

\item \stfhPropertyLabel{（零性）}若 $\Lambda(t)\cap\Lambda'(t)$ 的维数恒定，则
$\mu(\Lambda,\Lambda')=0$。

\item \stfhPropertyLabel{（直和性）}若 $S=S_1\oplus S_2$，则
\[
 \mu(\Lambda_1\oplus\Lambda_2,\Lambda'_1\oplus\Lambda'_2)
 =\mu(\Lambda_1,\Lambda'_1)+\mu(\Lambda_2,\Lambda'_2).
\]

\item \stfhPropertyLabel{（拼接性）}若 $a<c<b$，则
\[
 \mu(\Lambda,\Lambda')
 =\mu(\Lambda|_{[a,c]},\Lambda'|_{[a,c]})
 +\mu(\Lambda|_{[c,b]},\Lambda'|_{[c,b]}).
\]

\item \stfhPropertyLabel{（局部化）}若
$(S,\Omega)=(\mathbb R^{2n},\omega_0)$、
$\Lambda'(t)\equiv\mathbb R^n\times0$，且
$\Lambda(t)=\operatorname{Gr}(A(t))$，其中
$A\colon[a,b]\to\mathbb R^{n\times n}$ 是一条对称矩阵道路，则
\[
 \mu(\Lambda,\Lambda')
 =\frac12\operatorname{sign}A(b)-\frac12\operatorname{sign}A(a).
\]
\end{enumerate}

另一个重要而有用的性质，是它与定理 18.4.2 中康利--曾德尔指标的下述关系。为说明
这一点，还要考虑乘积辛向量空间
\[
 (S\times S,(-\Omega)\oplus\Omega),
\]
并把 $\operatorname{Gr}(\Psi)$ 和 $V\times V$ 看作其中的拉格朗日子空间。

\begin{theorem}
\label{thm:app-c-graph-maslov}
设 $S=\mathbb R^{2n}\cong\mathbb C^n$，
$\Psi\colon[a,b]\to\operatorname{Sp}(2n)$ 是辛道路，并令
$V=0\times\mathbb R^n$。则：
\begin{enumerate}[label=(\arabic*)]
\item 考虑对子 $(\Psi\cdot V,V)$，其中 $\Psi V$ 是
$t\mapsto\Psi(t)V$ 给出的拉格朗日道路。则
\[
 \mu(\Psi V,V)=\mu(\operatorname{Gr}(\Psi),V\times V).
\]
\item 若 $\Psi$ 如定理 18.4.2，即满足
$\Psi(a)=\mathbf1$ 且 $\Psi(b)V\cap V=0$，则
\[
 \mu(\Psi V,V)+\frac n2\in\mathbb Z.
\]
\end{enumerate}
\end{theorem}

(RS93) 中的\personRobbin{}与\personSalamon{}给出了定理 18.4.2 中康利--曾德尔指标的如下推广。

\begin{definition}
\label{def:app-c-generalized-cz}
设 $\Psi\colon[a,b]\to\operatorname{Sp}(2n)$ 是辛道路。定义 $\Psi$ 的康利--曾德尔
指标为
\[
 \mu_{\mathrm{CZ}}(\Psi)
 =\mu(\operatorname{Gr}(\Psi),\Delta)
 \in\mathbb Z+\frac12.
\]
\end{definition}

(RS93) 证明：当 $\Psi$ 满足 $\Psi(a)=\mathbf1$ 且
$\Psi(b)V\cap V=0$ 时，半整数 $\mu_{\mathrm{CZ}}(\Psi)$ 成为整数，并且与
定理 18.4.2 中的指标一致。
